
\begingroup\scriptsize\setlength{\tabcolsep}{2pt}
\begin{longtable}{@{}p{0.15\textwidth}p{0.21\textwidth}p{0.11\textwidth}p{0.13\textwidth}p{0.17\textwidth}p{0.07\textwidth}@{}}
\caption{Final scientific claims supported by the simulation campaign.}\label{tab:final_claims}\\
\toprule
\textbf{Claim} & \textbf{Supporting result} & \textbf{Fig.} & \textbf{Assumption} & \textbf{Limitation} & \textbf{Conf.}\\
\midrule
\endfirsthead
\toprule
\textbf{Claim} & \textbf{Supporting result} & \textbf{Fig.} & \textbf{Assumption} & \textbf{Limitation} & \textbf{Conf.}\\
\midrule
\endhead
HER2-antibody binding follows bounded saturation behavior across the modeled concentration range. & M1 occupancy increases monotonically with HER2 concentration and remains within [0, 1]. & Fig.~\ref{fig:m1_occupancy} & Reversible Langmuir-type binding with baseline $K_d=10$ pM. & No wet-lab binding calibration was performed. & High\\
Directional lymph-node-inspired flow changes sensor exposure compared with diffusion-only transport. & At $C=10$ pM, $v_{in}=5\times10^{-7}$ m/s gives 8.73 pM at 6000 s versus 4.98 pM for diffusion-only; time-to-50\% exposure decreases from 6038 s to 1839 s. & Fig.~\ref{fig:exp_a_flow} & Prescribed velocity is coupled into TDS convection. & M2B is a partially anatomical prescribed-velocity convection-diffusion extension, not a full anatomical lymph-node model and not a validated Darcy-flow pressure solution. & Medium\\
Sensor placement affects response timing and current response. & Subcapsular/cortical placement is fastest and strongest under reference flow; late-time current range is 559.9--600.5 pA under flow and 348.2--533.3 pA under diffusion-only transport. & Figs.~\ref{fig:exp_b_exposure}--\ref{fig:exp_b_current} & Three placement regions are post-processed over the same transport model. & Placement sweep is not a full remeshed COMSOL sensor relocation study. & Medium\\
Sub-pM or low-pM detectability is predicted only under favorable coupling, low-noise, and binding assumptions. & Experiment C gives 99 detectable rows and 36 failing rows across concentration, $\alpha_0$, noise, and $K_d$ sweeps. & Fig.~\ref{fig:exp_c_detectable} & Analytical GFET response with $\Delta I_{DS}\geq3\sigma$. & Graphene electronics and surface chemistry are simplified. & Medium\\
PBS-like Debye screening strongly reduces effective GFET signal for long receptor distances. & For PBS-like $\lambda_D=0.8$ nm and $d\geq5$ nm, 0 of 60 screened cases remain detectable. & Figs.~\ref{fig:exp_d_lod}--\ref{fig:exp_d_detectable} & Screened coupling model. & Debye lengths and distances are approximate feasibility assumptions. & Medium\\
Direct in-body HER2 detection likely requires short receptor distance, low-noise electronics, stronger transduction, or local screening mitigation. & Transport and placement improve antigen exposure, but detection fails under weak coupling, high noise, poor affinity, or PBS-like screening at longer distances. & Table~\ref{tab:design_requirements} & Synthesis of M2B exposure, M3 binding, M4 response, and Debye attenuation. & Simulation-based design conclusion, not clinical validation. & High\\
Feasible GFET-HER2 detection requires explicit design conditions on receptor distance, coupling, screening, and electronics noise. & Experiment E maps required $\alpha_0$, maximum allowed receptor-channel distance, and maximum allowed noise floor. For $K_d=10$ pM, 10 pA noise, and PBS-like screening at 1 nm, required $\alpha_0$ is 0.0326 at 1 pM and 0.00561 at 10 pM. & Fig.~\ref{fig:exp_e_requirement_map} & M3/M4 response with screened coupling. & Simulation-based requirement analysis, not fabricated-device performance. & Medium\\
\bottomrule
\end{longtable}
\endgroup
